\section{Introduction}
\label{sec:intro}

Training a dialogue policy against an LLM judge involves two choices that the literature treats
as separable. One is the \textbf{optimizer}: preference-pair methods in the DPO line
\citep{rafailov2023dpo}, which keep only the best and worst of the sampled candidates, versus
group-relative policy optimization \citep{shao2024deepseekmath}, which standardises the reward
across all sampled siblings and trains on every one. The other is the \textbf{reward horizon}:
whether the judge scores the candidate turn as produced, or the candidate plus a simulated
continuation of the dialogue --- the $K$-turn look-ahead introduced for preference-tree
optimisation in motivational interviewing by \citet{baruch2025pto}. Papers that vary one hold
the other fixed, and results are reported as if a horizon that helps one optimizer is a property
of the horizon.

We test that assumption directly, with a controlled $2 \times 2$: two optimizers $\times$ two
horizons, all four arms trained from the same 1B base model against the same simulated patient
population, with the same oracle, the same number of candidates per prompt, the same
temperatures, the same minimum-context rule, and the same number of iterations, and every model
state evaluated on eight instruments by two graders --- the training oracle, and a held-out judge
from a different model family that never touched training.

The assumption fails. \textbf{The ranking of the optimizers flips with the horizon.} With a
turn-level reward, the preference-tree method beats GRPO at the matched endpoint under both
graders ($\dz\,0.73$ / $1.27$); with a five-turn look-ahead, GRPO beats the preference-tree
method under both graders ($\dz\,-0.36$ / $-0.31$). Equivalently, by lever: look-ahead more than
doubles GRPO's gain over base, and is null-to-negative on the preference-tree method's outcome
scores in this regime. The interaction --- the difference-in-differences --- is
$\dz\,0.79$--$0.97$ depending on grader, larger than most published ablations in this space, and
every cell of it survives an independent re-draw of the contested endpoints' evaluation sets.

Three further results give the interaction its practical meaning. \textbf{Budget reverses the
rankings again} (\S\ref{sec:cost}): the four arms cost $8.1$ to $51.2$ reconstructed GPU-hours
for the same ten iterations, and at matched spend the cheapest arm (preference trees, $\Kz$) is
unbeaten until deep into the budget axis under the training oracle --- and on the held-out judge
is never beaten at all: it matches the most expensive arm's final gain ($+1.036$ vs.\ $+1.038$
rubric points) at $6.3\times$ less compute, a paired null that replicates across two evaluation
draws. \textbf{The horizon decides what is
learned, in both optimizers} (\S\ref{sec:behaviour}): under turn-level reward both $\Kz$ policies
drift into unearned praise --- a deterministic lexical marker, no model in the loop, finds
over-praise in $67\%$ (GRPO) and $21\%$ (PTO) of endpoint therapist turns --- and both $\Kf$
policies hold near base rates; in the preference-tree arm part of the suppressed inconsistency
relocates to advice-without-permission, in GRPO it does not. \textbf{The interaction is
regime-scoped} (\S\ref{sec:regime}): re-scoring the conversations of an earlier, easier regime of
the same task --- a 7B therapist, more cooperative patients, non-iterative training --- shows
look-ahead helping the preference-tree method there under two graders, so the null in the present
regime is a property of the regime, not of the optimizer pairing per se.

Our contributions: (i) the first controlled optimizer $\times$ reward-horizon factorial for
LLM-judged dialogue training, showing a sign-flipping interaction of $\dz \approx 0.8$--$1.0$
confirmed by two graders and replicated evaluation draws; (ii) a budget re-analysis on
reconstructed GPU-hours showing the practical optimizer ranking depends on spend, with the
cheapest arm unbeaten at low budgets; (iii) evidence from a judge-free transcript statistic that
the reward horizon, not the optimizer, determines whether the policy learns to flatter --- and
that the optimizer determines where the residual inconsistency goes; and (iv) a demonstration
that the same lever's sign changes across task regimes, which cautions against reporting reward
levers per optimizer or per benchmark as if they were free-standing.
